% !TeX root = ../../Book2.tex
% 纲要标题：### 11.3 幂等元与分解
% 纲要位置：计划纲要.md，第 1420 行。

\section{幂等元与分解}
\label{b2:sec:1103}

一个投影满足连续作用两次与作用一次相同，即 \(p^2=p\)。在含幺代数 \(A\) 内，若 \(e^2=e\)，右乘映射 \(R_e(a)=ae\) 就满足 \(R_e^2=R_e\)，并将正则左模投到左理想 \(Ae\)。环中的幂等元由此产生模的分解，但未必分离整个代数的乘法；区别来自这个幂等元是否位于中心。

% 纲要要点（供目录核对）：
%
% * 正交幂等元的模分解与中心幂等元的代数分解
% * 角环、Peirce块及其双模结构
% * 右乘投影、角块与模同态复合
%

\subsection{正交幂等元及单位分解}
\label{b2:subsec:1103:01}

\begin{definition}\label{b2:def:orthogonal-idempotents}
设 \(k\) 为域、\(A\) 为含幺结合 \(k\) 代数。若 \(e\in A\) 满足 \(e^2=e\)，则称 \(e\) 为\term[mideng-yuan]{幂等元}[idempotent]。若两个幂等元 \(e,f\in A\) 满足 \(ef=fe=0\)，则称它们正交。若有限族 \(e_1,\ldots,e_r\in A\) 两两正交且 \(\sum_i e_i=1_A\)，则称 \(1_A=e_1+\cdots+e_r\) 为单位元的正交幂等元分解。
\end{definition}

\begin{proposition}\label{b2:prop:regular-orthogonal-decomposition}
设 \(k\) 为域、\(A\) 为含幺结合 \(k\) 代数，\(1_A=e_1+\cdots+e_r\) 是正交幂等元分解。它给出正则模的分解
\[
{}_AA=\bigoplus_i Ae_i,\qquad A_A=\bigoplus_i e_iA.
\]
若每个 \(e_i\) 还位于中心，则它们给出代数同构 \(A\cong\prod_i Ae_i\)，其中因子 \(Ae_i\) 的单位为 \(e_i\)。
\end{proposition}
\begin{proof}
对任意 \(a\)，有 \(a=\sum_i ae_i\)。若 \(\sum_i a_i e_i=0\)，右乘 \(e_j\) 只留下 \(a_j e_j\)，所以各项为零，左模分解成立。左乘的同样计算给出右模分解。

当各 \(e_i\) 在中心时，\(Ae_i\) 是双边理想和有单位的子代数，不同因子的乘积为零。映射 \(a\mapsto(ae_i)_i\) 保持加法、标量与乘法，并把单位送到各因子的单位族。其逆为求和映射，由刚证的直和分解互逆，得到代数同构。
\end{proof}

非中心幂等元只保证模分解。例如 \(A=M_2(k)=AE_{11}\oplus AE_{22}\) 是按列的左模分解，但 \(E_{11}\in AE_{11}\)、\(E_{12}\in AE_{22}\)，且 \(E_{11}E_{12}=E_{12}\ne0\)。这两个左模因子之间仍有非零乘积，不能成为直积代数中互相乘得零的两个因子。

\subsection{角环与矩阵分块}
\label{b2:subsec:1103:02}

设 \(k\) 为域、\(A\) 为含幺 \(k\) 代数、\(e\in A\) 为幂等元。集合 \(eAe\) 对加法、标量及乘法封闭，且以 \(e\) 为单位元；它称为 \(A\) 对应于 \(e\) 的\term[jiao-huan]{角环}[corner ring]，在当前情形也可称角代数。若 \(e\ne1_A\)，包含 \(eAe\subseteq A\) 不保持单位元，因而不是本书约定的含幺代数同态。

\ProseParagraphBreak
在 \(A=M_n(k)\) 中，若 \(e=\operatorname{diag}(I_s,0)\)、\(1\le s\le n\)，两侧乘 \(e\) 就只保留左上 \(s\times s\) 块，其余各块为零，所以 \(eAe\cong M_s(k)\)。对单位元的正交分解 \(1=\sum_i e_i\)，\(e_i a e_j\) 同样选出第 \((i,j)\) 块；下面的分解说明这种矩阵分块计算在一般代数中也成立。

\begin{proposition}[幂等元分块]\label{b2:prop:peirce-decomposition}
设 \(k\) 为域，\(A\) 为含幺结合 \(k\) 代数，\(1_A=e_1+\cdots+e_r\) 为正交幂等元分解，则
\begin{equation}\label{b2:eq:peirce-decomposition}
A=\bigoplus_{i,j}e_iAe_j
\end{equation}
作为向量空间成立。乘法满足 \((e_iAe_j)(e_sAe_t)=0\) 当 \(j\ne s\)；当 \(j=s\) 时，乘积属于 \(e_iAe_t\)。
每个块 \(e_iAe_j\) 通过 \(A\) 内的乘法成为 \((e_iAe_i,e_jAe_j)\) 双模，左右单位分别为 \(e_i\) 与 \(e_j\)。
\end{proposition}
\begin{proof}
把 \(a=1a1\) 展开，得到 \(a=\sum_{i,j}e_i a e_j\)。若这类分量之和为零，左乘 \(e_s\)、右乘 \(e_t\) 就只留下第 \((s,t)\) 个分量，所以分解唯一。两分量相乘时中间出现 \(e_je_s\)，它在不同指标时为零，相同指标时为 \(e_j\)，给出所述乘法规则。角环 \(e_iAe_i\) 左乘及 \(e_jAe_j\) 右乘都保持 \(e_iAe_j\)，并由结合律满足左右作用的相容性；对其中的 \(x\)，有 \(e_ix=x=xe_j\)，所以这些作用幺性。
\end{proof}

这种分块称为\term[Peirce-fenjie]{Peirce 分解}[Peirce decomposition]。\footnote{本杰明·皮尔斯（Benjamin Peirce，1809-1880），美国数学家、天文学家，在《Linear Associative Algebra》中以幂等元研究结合代数的结构。}因此可以把 \(a\) 记作分块矩阵 \((e_i a e_j)\)。乘积第 \((i,j)\) 块为 \(\sum_s(e_i a e_s)(e_s b e_j)\)，正是矩阵乘法式。不同块可能是不同的向量空间，不能未经论证就把它们全当作同一个系数环；非对角块作为双模连接两端的角环。

\ProseParagraphBreak
在 \(M_n(k)\) 中取 \(e_i=E_{ii}\)，各角环为 \(kE_{ii}\)，各块为 \(kE_{ij}\)，于是恢复通常的矩阵条目。若改取上三角代数 \(T_2(k)\)，右上块 \(e_1Ae_2=kE_{12}\) 非零，左下块 \(e_2Ae_1=0\)。两个角环都同构于 \(k\)，但整个代数不是 \(k\times k\)：非零右上块仍参与乘法。

\subsection{模分解与环中幂等元}
\label{b2:subsec:1103:03}

前面模论中已经见过投影的像与核分解；在\cref{b2:prop:finite-projective-summand}中，有限生成投射模还由有限秩自由模上的幂等自同态刻画。这里重新写出任意模的投影分解，是为了与环中幂等元产生的右乘投影对应起来。

\begin{proposition}\label{b2:prop:idempotent-module-decomposition}
对任意含幺环 \(R\) 和幺性左 \(R\) 模 \(M\)，满足 \(p^2=p\) 的 \(p\in\operatorname{End}_R(M)\) 给出
\[
M=\operatorname{im}p\oplus\ker p.
\]
反过来，每个指定的直和分解都给出投影幂等元。一组有限两两正交且和为恒等的投影，对应有限个子模的内直和分解。
\end{proposition}
\begin{proof}
写 \(m=p(m)+\bigl(m-p(m)\bigr)\)，第二项在核中；若 \(x=p(y)\) 又满足 \(p(x)=0\)，则 \(x=p^2(y)=p(x)=0\)，所以和为直和。反向取沿其余分量投到指定分量的映射，便满足平方为自身。有限族时，和为恒等保证生成，两两正交保证从一个分量到另一个的投影为零，从而分解唯一。
\end{proof}

在正则左模上，\cref{b2:prop:regular-bimodule-endomorphism}说明模自同态是右乘，所以幂等元 \(e\) 给出的投影为 \(a\mapsto ae\)，像为 \(Ae\)，核为 \(A(1-e)\)。左乘 \(L_e(a)=ea\) 虽也满足 \(L_e^2=L_e\)，却未必是左模自同态：左线性要求 \(e(ra)=r(ea)\)，取 \(a=1\) 就得到 \(er=re\) 对所有 \(r\in A\) 成立。反过来，\(e\) 在中心时此条件确实成立。取 \(A=M_2(k)\)、\(e=E_{11}\)、\(r=E_{21}\)，则
\[
L_e(r)=E_{11}E_{21}=0,
\qquad rL_e(1)=E_{21}E_{11}=E_{21}\ne0.
\]
因此左正则模中的投影须用右乘；不能仅凭算子平方等于自身，就把它当作左模投影。

\begin{proposition}\label{b2:prop:corner-endomorphism}
设 \(A\) 为含幺环，\(e,f\in A\) 为幂等元。取值映射给出交换群同构
\[
\operatorname{Hom}_A(Ae,Af)\longrightarrow eAf,
\qquad h\longmapsto h(e).
\]
它的逆把 \(b\in eAf\) 送到 \(h_b:Ae\rightarrow Af\)、\(h_b(ae)=ab\)。若 \(g\in A\) 也是幂等元，\(b\in eAf\)、\(c\in fAg\)，则
\[
h_c\circ h_b=h_{bc}:Ae\longrightarrow Ag.
\]
特别地，有含幺环同构
\[
(eAe)^{\mathrm{op}}\longrightarrow\operatorname{End}_A(Ae),
\qquad b^{\mathrm{op}}\longmapsto h_b.
\]
若 \(k\) 为域、\(A\) 为含幺结合 \(k\) 代数，则上述取值同构为 \(k\) 线性同构，上述自同态环同构为含幺 \(k\) 代数同构。
\end{proposition}
\begin{proof}
若 \(h:Ae\to Af\) 左线性，因为 \(h(e)\in Af\)，可写成 \(h(e)=a_0f\)。又 \(h(e)=h(e^2)=e h(e)\)，所以
\[
h(e)=ea_0f\in eAf.
\]
对任意 \(a\in A\)，左线性又迫使 \(h(ae)=a\cdot h(e)\)，故 \(e\) 的像决定整个同态。

反过来，取 \(b\in eAf\)。若 \(ae=a'e\)，则 \((a-a')e=0\)，而 \(eb=b\)，所以 \((a-a')b=(a-a')eb=0\)，证明 \(h_b(ae)=ab\) 与代表元选择无关。由于 \(bf=b\)，有 \(ab\in Af\)；结合律给出左 \(A\) 线性。取值映射与 \(b\mapsto h_b\) 都保持加法，而 \(h_b(e)=eb=b\)，所以这两个加法群映射互为逆。

对 \(b\in eAf\)、\(c\in fAg\)，有 \(bc\in eAg\)，并且
\[
(h_c\circ h_b)(ae)=h_c(ab)=abc=h_{bc}(ae).
\]
这里 \(ab=(ab)f\) 作为 \(Af\) 中的元素，正可代入 \(h_c\) 的定义。取 \(f=g=e\)，便有 \(h_b\circ h_c=h_{cb}\)，所以参数乘法相对于通常的自同态复合恰好反转，得到 \((eAe)^{\mathrm{op}}\) 到自同态环的同构；\(h_e(ae)=ae\)，故角环的单位 \(e\) 对应恒等映射。当 \(A\) 还是 \(k\) 代数时，取值与其逆都保持 \(k\) 标量，故这些对应分别是所述的 \(k\) 线性同构与 \(k\) 代数同构。
\end{proof}

这个取值对应说明，Peirce 块 \(eAf\) 同时记录从左理想 \(Ae\) 到 \(Af\) 的所有左模同态，块乘法 \((eAf)(fAg)\subseteq eAg\) 对应先施加 \(h_b\)、再施加 \(h_c\) 的复合。取 \(f=e\)，便把正则模的分解细化到某个直和因子内部。后续判别 \(Ae\) 是否还能分解时，考察的是它的自同态环中有无非平凡幂等元，也就是角环 \(eAe\) 中的同一问题。
